% TEMPLATE-MILC-v3.tex - plantilla maestra para todas las guias MILC v3
% Ver scripts/generadores/build-guia-milc.py para la lista completa de
% claves esperadas. El script valida que no queden placeholders sin
% reemplazar antes de escribir el archivo final.

\documentclass[11pt,letterpaper]{article}
\usepackage[margin=1.75cm]{geometry}
\usepackage[table]{xcolor}
\usepackage{fontspec}
\usepackage[spanish]{babel}
\usepackage{tikz}
\usepackage[most]{tcolorbox}
\usepackage{tabularx}
\usepackage{array}
\usepackage{fancyhdr}
\usepackage{titlesec}
\usepackage{ragged2e}
\usepackage{enumitem}
\usepackage{microtype}

\setmainfont{Helvetica Neue}
\setsansfont{Helvetica Neue}
\setlength{\parindent}{0pt}
\setlength{\parskip}{6pt}
\hyphenpenalty=200
\exhyphenpenalty=200
\pretolerance=400
\tolerance=2000
\setlength{\emergencystretch}{1em}
\renewcommand{\arraystretch}{1.25}
\setlist{leftmargin=5mm,itemsep=1mm,topsep=1mm}
\newcolumntype{Y}{>{\RaggedRight\arraybackslash}X}

\definecolor{milcMagenta}{HTML}{E6007E}
\definecolor{milcVino}{HTML}{5A0038}
\definecolor{milcTurquesa}{HTML}{0093A5}
\definecolor{milcVerde}{HTML}{58A923}
\definecolor{milcMostaza}{HTML}{E5B400}
\definecolor{milcOcre}{HTML}{B86600}
\definecolor{milcNegro}{HTML}{241816}
\definecolor{milcGris}{HTML}{F7F7F4}
\definecolor{milcLinea}{HTML}{E8E2DD}
\definecolor{milcGradePrimary}{HTML}{0066FF}
\definecolor{milcGradeDark}{HTML}{003D99}
\definecolor{milcGradeSoft}{HTML}{D6E8FF}
\definecolor{milcGradeAccent}{HTML}{CFE7FF}
\definecolor{milcGradeText}{HTML}{FFFFFF}
\color{milcNegro}

\pagestyle{fancy}
\fancyhf{}
\lhead{\small Tecnología e Informática · Grado 11}
\rhead{\small Guía 14 · MILC v3}
\cfoot{\small\thepage}
\renewcommand{\headrulewidth}{0pt}

\newtcolorbox{titlebox}[1]{
  enhanced,colback=#1,colframe=#1,arc=7mm,boxrule=0pt,
  left=7mm,right=7mm,top=3mm,bottom=3mm,
  before skip=8pt,after skip=8pt,
  fontupper=\sffamily\bfseries\Large\color{white}
}

\newtcolorbox{softbox}[3]{
  enhanced,colback=#2,colframe=#1,borderline west={4pt}{0pt}{#1},
  arc=2mm,boxrule=.4pt,left=4mm,right=4mm,top=3mm,bottom=3mm,
  before skip=6pt,after skip=8pt,title={#3},coltitle=white,
  colbacktitle=#1,fonttitle=\sffamily\bfseries,fontupper=\RaggedRight,
  attach boxed title to top left={xshift=3mm,yshift=-2mm},
  boxed title style={arc=2mm, boxrule=0pt}
}

\newtcolorbox{drawbox}[2]{
  enhanced,colback=white,colframe=#1,arc=4mm,boxrule=1pt,
  left=5mm,right=5mm,top=4mm,bottom=4mm,before skip=8pt,after skip=8pt,
  title={#2},coltitle=white,colbacktitle=#1,fonttitle=\sffamily\bfseries,
  fontupper=\RaggedRight,
  attach boxed title to top left={xshift=4mm,yshift=-2mm},
  boxed title style={arc=3mm, boxrule=0pt}
}

\newtcolorbox{infoband}[1]{
  enhanced,colback=#1!8,colframe=#1!50,arc=2mm,boxrule=.5pt,
  left=4mm,right=4mm,top=2mm,bottom=2mm,before skip=4pt,after skip=4pt,
  fontupper=\small\RaggedRight
}

% Puente narrativo entre secciones (contrato editorial Conéctate).
% Da continuidad: "Acabas de X. Ahora vas a Y porque Z."
% Visualmente: banda izquierda en color del grado, texto en itálica pequeña.
\newtcolorbox{puentebox}{
  enhanced,colback=milcGradeSoft,colframe=milcGradeDark,
  borderline west={3pt}{0pt}{milcGradeDark},
  arc=0mm,boxrule=0pt,left=5mm,right=4mm,top=2.5mm,bottom=2.5mm,
  before skip=4pt,after skip=8pt,
  fontupper=\itshape\small\color{milcNegro}\RaggedRight
}
\newcommand{\puente}[1]{%
  \begin{puentebox}%
    \textcolor{milcGradeDark}{\textbf{\textrightarrow}}\hspace{2mm}#1%
  \end{puentebox}%
}

\newcommand{\linead}[1]{\noindent\color{milcLinea}\rule{#1}{0.5pt}\color{milcNegro}}
\newcommand{\respuesta}[1]{\par\vspace{2mm}\foreach \n in {1,...,#1}{\noindent\color{milcLinea}\rule{\linewidth}{0.5pt}\par\vspace{5mm}}\color{milcNegro}}
\newcommand{\checkbox}{\textcolor{milcGradeDark}{\fbox{\phantom{x}}}\hspace{5pt}}

\newcommand{\phasecard}[4]{
\begin{tcolorbox}[enhanced,colback=#3,colframe=#2,arc=3mm,boxrule=.5pt,height=3.05cm,valign=center,halign=center,left=1.5mm,right=1.5mm,top=2mm,bottom=2mm]
{\sffamily\bfseries\fontsize{22}{24}\selectfont\textcolor{#2}{#1}}\par
{\sffamily\bfseries\small #4}
\end{tcolorbox}
}

\begin{document}
\thispagestyle{empty}
\begin{tikzpicture}[remember picture,overlay]
  \fill[white] (current page.south west) rectangle (current page.north east);
  \path[fill=milcGradePrimary,rounded corners=16mm]
    ([xshift=.55cm,yshift=-.55cm]current page.north west) rectangle
    ([xshift=-.55cm,yshift=3.65cm]current page.south east);

  \node[anchor=north west,text=milcGradeText] at ([xshift=2.0cm,yshift=-1.85cm]current page.north west)
    {\sffamily\bfseries\fontsize{14}{16}\selectfont GRADO};
  \node[anchor=north west,text=milcGradeText,opacity=.82] at ([xshift=2.0cm,yshift=-2.75cm]current page.north west)
    {\sffamily\bfseries\fontsize{9}{11}\selectfont SERIE GUÍAS · TECNOLOGÍA E INFORMÁTICA};

  \begin{scope}[shift={([xshift=-3.05cm,yshift=-2.55cm]current page.north east)}]
    \fill[white,opacity=.80] (-.62,-.52) rectangle (.62,.52);
    \draw[milcGradeText,line width=1pt] (-.62,-.52) rectangle (.62,.52);
    \fill[milcGradeDark] (-.42,-.38) rectangle (-.22,.06);
    \fill[milcGradeAccent] (-.10,-.38) rectangle (.10,.24);
    \fill[milcGradePrimary!60!black] (.22,-.38) rectangle (.42,.38);
  \end{scope}

  \node[anchor=west,text=milcGradeText] at ([xshift=1.9cm,yshift=-12.85cm]current page.north west)
    {\sffamily\bfseries\fontsize{124}{124}\selectfont 11};
  \draw[milcGradeText,line width=1.7pt]
    ([xshift=4.52cm,yshift=-11.68cm]current page.north west) circle (.13cm);

  \node[anchor=west,text=milcGradeText,text width=8.7cm,align=left] at ([xshift=10.35cm,yshift=-11.6cm]current page.north west)
    {
     {\sffamily\bfseries\fontsize{19}{22}\selectfont Hojas de cálculo ---\\como\\base de datos}\\[4mm]
     {\sffamily\bfseries\fontsize{23}{26}\selectfont Guía 14-11-TIC}\\[2mm]
     {\sffamily\fontsize{13}{16}\selectfont Período 2 · Automatización y procesos}\\[5mm]
     {\sffamily\bfseries\fontsize{11}{13}\selectfont Institución Educativa Sor María Juliana}\\[1mm]
     {\sffamily\fontsize{10}{12}\selectfont Autor: Dr. Álvaro Cárdenas Orozco}
    };

  \path[fill=milcGradeSoft,rounded corners=7mm]
    ([xshift=1.35cm,yshift=.85cm]current page.south west) rectangle
    ([xshift=-1.35cm,yshift=4.25cm]current page.south east);
  \draw[milcGradeDark,line width=.65pt,rounded corners=7mm]
    ([xshift=1.35cm,yshift=.85cm]current page.south west) rectangle
    ([xshift=-1.35cm,yshift=4.25cm]current page.south east);
  \node[anchor=center,text=milcNegro,text width=17.0cm,align=center] at ([yshift=2.55cm]current page.south)
    {\sffamily\fontsize{10}{12}\selectfont Metodología MILC · Apertura ancestral · Pensamiento computacional · Triángulo Dussel-Estoicismo-Floridi\\
    \textcolor{milcGradeDark}{\bfseries Producto:} 1 hoja de cálculo estructurada con datos reales (mínimo 25 filas) configurada con 3 vistas funcionales (filtro avanzado por categoría, ordenamiento por columna numérica, fórmulas que calculan métricas clave) lista para usarse como base de datos operativa de un proceso};
\end{tikzpicture}
\mbox{}
\newpage

\section*{Apertura: Hojas de cálculo como base de datos}

\begin{softbox}{milcGradeDark}{milcGradeSoft}{Saber ancestral}
El \textbf{libro contable del mercado campesino} de Cartago es la primera base de datos: cada fila una venta, cada columna un dato (\emph{qué}, \emph{cuánto}, \emph{a quién}, \emph{por cuánto}, \emph{cómo pagó}). El \textbf{libro mayor del tendero} del barrio cumple la misma función: registro estructurado, una observación por fila, una variable por columna. La \textbf{lechera} llevaba su cuaderno con tantas filas como vacas y tantas columnas como días del mes. Excel y Google Sheets son la evolución natural del libro mayor: misma anatomía (fila = caso, columna = variable), misma disciplina (un dato en su celda exacta), pero con cálculo automático, ordenamiento instantáneo y filtros sin reescribir nada.
\end{softbox}

\begin{softbox}{milcTurquesa}{milcGris}{Saber contemporáneo}
Una \textbf{hoja como base de datos} cumple 3 reglas innegociables: (1) \textbf{una observación por fila}; (2) \textbf{una variable por columna}; (3) \textbf{encabezados claros en la primera fila}. Una vez cumplidas, la hoja admite 3 vistas operativas: (a) \textbf{filtro} para aislar subconjuntos; (b) \textbf{orden} para priorizar; (c) \textbf{fórmulas} para calcular métricas (\texttt{SUMA}, \texttt{PROMEDIO}, \texttt{CONTAR.SI}, \texttt{BUSCARV}). Esa estructura mínima --- llamada \emph{forma tabular limpia} (\emph{tidy data}) --- es la diferencia entre una hoja que sirve como base de datos y un montón de números en celdas.
\end{softbox}

\begin{softbox}{milcMostaza}{milcGris}{Pregunta-puente}
\textit{¿Cómo sabía el tendero qué columna agregar a su libro mayor antes de tener un curso de Excel? ¿Y qué pierde un emprendedor novato cuando construye una hoja con datos mezclados en celdas (\emph{``Juan, 5 panes, jueves''} en una sola casilla) sin separar las variables?}
\end{softbox}

\begin{softbox}{milcVerde}{milcGris}{Saber y saber hacer}
Al terminar podrás: (1) \textbf{identificar} las 3 reglas innegociables de una hoja como base de datos; (2) \textbf{explicar} con tus palabras la diferencia entre datos limpios (\emph{tidy}) y datos sucios; (3) \textbf{aplicar} las 3 vistas (filtro, orden, fórmulas) sobre una hoja real con datos del estudiante; (4) \textbf{evaluar} si la hoja sirve como base de datos operativa o si necesita rediseño antes de automatizar lo que va encima.
\end{softbox}

\small
\begin{tabularx}{\textwidth}{>{\bfseries}p{3.0cm}Y}
\rowcolor{milcGradePrimary!12}Autor & Dr. Álvaro Cárdenas Orozco\\
DBA & Diseño y mejora de procesos digitales con criterio ético (MEN, T\&I)\\
Referentes & Pensamiento computacional · Lean / Kaizen · Floridi (ética informacional)\\
Producto final & 1 hoja de cálculo estructurada con datos reales (mínimo 25 filas) configurada con 3 vistas funcionales (filtro avanzado por categoría, ordenamiento por columna numérica, fórmulas que calculan métricas clave) lista para usarse como base de datos operativa de un proceso\\
\end{tabularx}
\normalsize

\puente{Antes de empezar, mira el viaje completo. En la sesión 3 capturaste información con un formulario. Hoy aprendes a \emph{operar} la hoja que recibe esa información para que sirva como base de datos real: filtrar, ordenar, calcular. Sin esta sesión, la hoja que llena el formulario es solo \emph{una pila de datos}; con esta sesión, es un sistema operativo de información.}

\begin{titlebox}{milcGradeDark}
Ruta MILC: así vamos a aprender
\end{titlebox}
\begin{tabularx}{\textwidth}{>{\centering\arraybackslash}X>{\centering\arraybackslash}X>{\centering\arraybackslash}X>{\centering\arraybackslash}X}
\phasecard{1}{milcVerde}{milcGris}{Escuta\\[-1mm]\small Observo, escucho y pregunto.} &
\phasecard{2}{milcTurquesa}{milcGris}{Sistematización\\[-1mm]\small Pensamiento computacional.} &
\phasecard{3}{milcMagenta}{milcGris}{Praxis\\[-1mm]\small Construyo y aplico.} &
\phasecard{4}{milcVino}{milcGris}{Evaluación\\[-1mm]\small Reflexión liberadora.}
\end{tabularx}


\puente{Antes de teorizar, vas a abrir 2 hojas reales: una bien estructurada y otra desordenada (datos mezclados, columnas duplicadas, celdas combinadas). Vas a intentar responder 3 preguntas con cada una y a comparar cuánto tardas. La hoja sucia es lección por contraste.}

\begin{titlebox}{milcVerde}
Fase 1 · Escuta
\end{titlebox}

\begin{softbox}{milcVerde}{milcGris}{Escena para escuchar}
\textbf{Actividad 1 · IDENTIFICA --- Hoja limpia vs hoja sucia} (15 min · individual). El docente comparte 2 hojas reales: una limpia (25 filas, 5 columnas bien tipadas) y una sucia (datos mezclados, celdas combinadas, columnas con texto y número juntos). Sobre cada una intenta responder en 5 minutos: \emph{(1) ¿cuántas filas cumplen una condición específica?}, \emph{(2) ¿cuál es el promedio de una columna numérica?}, \emph{(3) ¿cuál fila tiene el valor máximo?}. Cronometra el tiempo en cada hoja y anota qué te frustró en la sucia.
\end{softbox}

\checkbox Intenté responder las 3 preguntas en la hoja limpia y cronometré.\\
\checkbox Intenté responder las 3 preguntas en la hoja sucia y cronometré.\\
\checkbox Anoté qué problemas concretos encontré en la hoja sucia (columnas, celdas, formatos).

\begin{infoband}{milcVerde}
\textbf{Lo que va al cuaderno (Actividad 1 · IDENTIFICA).} Título: «Actividad 1 --- Hoja limpia vs hoja sucia». \textbf{Formato:} tabla 2 filas (limpia / sucia) × 3 columnas (Tiempo total~|~Frustración principal~|~Pregunta que no pude responder). \textbf{Sabes que terminaste cuando} puedes nombrar 2 problemas concretos de la hoja sucia con vocabulario técnico (\emph{``celdas combinadas''}, \emph{``columna con dos tipos''}, etc.).
\end{infoband}

\puente{Ya viste el contraste. Ahora vamos a entender la \textbf{anatomía} de una hoja como base de datos: las 3 reglas innegociables, los errores típicos que rompen la estructura, las 3 vistas operativas que la hacen funcional.}

\begin{titlebox}{milcTurquesa}
Fase 2 · Sistematización con pensamiento computacional
\end{titlebox}

\begin{softbox}{milcTurquesa}{milcGris}{Cuatro pilares aplicados al tema}
Las 3 reglas innegociables y las 3 vistas operativas son el alfabeto de toda hoja profesional. Vamos a descomponerlas con los pilares del pensamiento computacional para que toda decisión de estructura tenga razón, no costumbre.
\end{softbox}

\small
\begin{tabularx}{\textwidth}{>{\bfseries\RaggedRight\arraybackslash}p{3.6cm}Y}
\rowcolor{milcTurquesa!90}\color{white}Pilar & \color{white}\bfseries Aplicación al tema\\
1. Descomposición & La estructura se descompone en 3 reglas: (1) \textbf{una fila = una observación}: cada fila representa un caso completo (una venta, un cliente, un día). (2) \textbf{una columna = una variable}: cada columna contiene un solo tipo de dato (texto, número, fecha) y no mezcla. (3) \textbf{encabezados claros en la primera fila}: nombres cortos sin espacios ni acentos (\emph{nombre\_cliente}, \emph{fecha\_pedido}, \emph{monto\_total}) facilitan fórmulas y conexión a otras herramientas.\\
2. Reconocimiento de patrones & Las 3 vistas operativas siguen el patrón \textbf{``preguntar a la base con criterio''}: (a) \textbf{Filtro}: \emph{¿cuáles filas cumplen X?} --- aplicado sobre 1 o más columnas; (b) \textbf{Orden}: \emph{¿cuál fila lidera o queda al final?} --- por columna numérica o fecha; (c) \textbf{Fórmulas}: \emph{¿cuánto es el total, el promedio, el conteo?} --- aplicado sobre rangos completos. Cada vista responde una pregunta distinta; combinarlas es la habilidad real.\\
3. Abstracción & Lo esencial cabe en una pregunta: \emph{¿la hoja responde a la pregunta sin reescribir nada?}. Si para responder hay que copiar, pegar, reestructurar o crear una pestaña nueva, la hoja no es base de datos: es un documento estático. La prueba operativa: \emph{aplica filtro + orden + fórmula y la respuesta aparece}.\\
4. Algoritmo conceptual & Algoritmo: (1) revisa las 3 reglas (fila = caso, columna = variable, encabezados); (2) define las 3-5 preguntas operativas más frecuentes; (3) para cada pregunta, decide qué vista la responde (filtro / orden / fórmula); (4) configura las 3 vistas; (5) pruébalas; (6) deja una pestaña \emph{``Cómo usar''} con 1 línea por vista.\\
\end{tabularx}
\normalsize

\begin{drawbox}{milcTurquesa}{Anatomía de una hoja como base de datos}
\textbf{Pestaña 1 (Datos):} 1 fila = 1 caso; 1 columna = 1 variable. \\[1mm] \textbf{Encabezados:} fila 1, sin espacios ni acentos. \\[1mm] \textbf{Tipos:} cada columna con tipo coherente (número, texto, fecha). \\[1mm] \textbf{Pestaña 2 (Vistas):} 3 vistas configuradas (filtro, orden, fórmulas) con 1 frase de uso. \\[1mm] \textbf{Pestaña 3 (Cómo usar):} 1 línea por vista que explique qué pregunta responde y a quién sirve.
\end{drawbox}

\begin{softbox}{milcMostaza}{milcGris}{Errores comunes y su corrección}
(1) Celdas combinadas: rompen el filtro y la fórmula. (2) Columna mezclada (\emph{``5 panes, jueves''} en una sola celda): no se puede ordenar ni calcular. (3) Encabezados con espacios o acentos: rompen integraciones (Apps Script, Zapier, fórmulas que llaman columnas por nombre). (4) Hojas con totales \emph{``al final''}: \texttt{=SUMA(A:A)} se rompe al agregar filas si el total queda incluido. (5) Una sola pestaña para todo: mezcla datos crudos, vistas y notas; pierde claridad operativa.
\end{softbox}

\begin{infoband}{milcTurquesa}
\textbf{Lo que va al cuaderno (Actividad 2 · EXPLICA).} Título: «Actividad 2 --- Anatomía de la hoja-base». \textbf{Formato:} ficha con las 3 reglas + 3 vistas + 5 errores comunes con marca 🚨 del que más temes cometer + 1 ejemplo personal por regla. \textbf{Sabes que terminaste cuando} podrías corregir una hoja ajena nombrando los errores con vocabulario técnico.
\end{infoband}

\puente{Tienes el método. Toca aplicarlo: una hoja real con 25+ filas, 3 vistas configuradas (filtro avanzado, orden, fórmulas clave), lista para servir como base de datos del proceso que estás automatizando desde la sesión 1.}

\begin{titlebox}{milcMagenta}
Fase 3 · Praxis con pensamiento computacional
\end{titlebox}

\begin{softbox}{milcMagenta}{milcGris}{Construyendo el producto con los cuatro pilares}
\textbf{Actividad 3 · APLICA --- Hoja con 3 vistas operativas} (30 min · individual). Hora de aplicar. Vas a tomar datos reales (puedes usar la hoja enlazada al formulario de sesión 3, los datos de un proceso real cercano, o un dataset descargado) y vas a configurarla como base de datos operativa.
\end{softbox}

\small
\begin{tabularx}{\textwidth}{>{\bfseries\RaggedRight\arraybackslash}p{3.6cm}Y}
\rowcolor{milcMagenta!90}\color{white}Pilar & \color{white}\bfseries Aplicación al producto\\
Descomposición & Tu hoja se descompone en 5 piezas: (1) pestaña \emph{Datos} con 25+ filas y encabezados limpios; (2) tipos coherentes por columna; (3) pestaña \emph{Vistas} con filtro avanzado configurado; (4) en \emph{Vistas}, orden por columna numérica o fecha; (5) en \emph{Vistas}, 3-5 fórmulas que calculen métricas clave (\texttt{SUMA}, \texttt{PROMEDIO}, \texttt{CONTAR.SI}, \texttt{BUSCARV}, \texttt{UNICOS}).\\
Patrones & Sigue el patrón \textbf{``cada vista responde 1 pregunta operativa real''}: antes de configurar la vista, escribe la pregunta. Si no hay pregunta, no hay vista. \emph{``Filtro por ciudad porque queremos ver solo Cartago''} sí; \emph{``filtro por ciudad porque sí''} no.\\
Abstracción & Cada fórmula expresa una sola decisión: \texttt{=PROMEDIO(D2:D26)} responde \emph{cuánto típicamente}; \texttt{=CONTAR.SI(C:C,``9A'')} responde \emph{cuántos casos de 9A hay}. No metas 3 fórmulas anidadas si no las necesitas: una fórmula clara vale por dos anidadas \emph{``elegantes''}.\\
Algoritmo de armado & Algoritmo: (1) abre la hoja; (2) verifica las 3 reglas; (3) limpia encabezados; (4) crea pestaña \emph{Vistas}; (5) configura filtro, orden y fórmulas; (6) crea pestaña \emph{Cómo usar} con 1 frase por vista; (7) comparte la URL con permisos de \emph{``solo lectura''} para alguien externo.\\
\end{tabularx}
\normalsize

\begin{drawbox}{milcMagenta}{Checklist antes de entregar la hoja}
\checkbox 3 reglas de tidy data cumplidas (fila = caso, columna = variable, encabezados limpios). \\[1mm] \checkbox 25+ filas con datos reales. \\[1mm] \checkbox Pestaña \emph{Vistas} con filtro, orden y fórmulas configurados. \\[1mm] \checkbox Cada vista responde una pregunta operativa real. \\[1mm] \checkbox Pestaña \emph{Cómo usar} con 1 frase por vista. \\[1mm] \checkbox URL compartible con permisos correctos.
\end{drawbox}

\begin{softbox}{milcMagenta}{milcGris}{Plantilla de trabajo}
\textbf{Hoja.} \textit{Tema + dueño + URL.} \\[1mm] \textbf{Datos.} \textit{25+ filas, encabezados.} \\[1mm] \textbf{Vista 1 (filtro).} \textit{Pregunta + columna + condición.} \\[1mm] \textbf{Vista 2 (orden).} \textit{Pregunta + columna + dirección.} \\[1mm] \textbf{Vista 3 (fórmulas).} \textit{3-5 fórmulas con su pregunta de uso.}
\end{softbox}

\begin{infoband}{milcMagenta}
\textbf{Lo que va al cuaderno (Actividad 3 · APLICA).} Título: «Actividad 3 --- Mi hoja como base». \textbf{Formato:} pantallazo de la pestaña \emph{Datos} + pantallazo de \emph{Vistas} + checklist marcado + URL pegada. \textbf{Extensión:} 1 página de cuaderno. \textbf{Sabes que terminaste cuando} un compañero abrió tu hoja sin tu explicación y supo responder las 3 preguntas operativas.
\end{infoband}

\puente{Lo que sigue resume \emph{qué} entregas hoy: 1 hoja operativa con 3 vistas + 1 frase por vista que explique \emph{qué decisión apoya}. El criterio principal: que otro estudiante, abriendo tu hoja, pueda entender en 2 minutos qué se puede preguntar a la base y dónde.}

\begin{titlebox}{milcMostaza}
Producto de la sesión
\end{titlebox}
\textbf{1 hoja operativa con 3 vistas configuradas + pestaña de uso}.
\begin{softbox}{milcMostaza}{milcGris}{Criterios de calidad}
(1) 3 reglas tidy cumplidas. (2) 25+ filas reales en pestaña Datos. (3) 3 vistas operativas configuradas. (4) Cada vista responde una pregunta operativa real. (5) Pestaña Cómo usar legible para externos.
\end{softbox}

\puente{Antes de cerrar, mira la hoja desde las cinco dimensiones humanas. Estructurar bien es respeto por quien va a consultar mañana; estructurar mal es regalarse trabajo extra y heredar caos.}

\begin{titlebox}{milcVino}
Fase 4 · Evaluación liberadora — cinco dimensiones
\end{titlebox}

\begin{softbox}{milcVino}{milcGris}{Sentido de la evaluación}
Evaluar no es solo revisar una nota. Es reconocer qué aprendí, qué cambió en mí, qué emociones manejé, cómo me relaciono con la ciudad y la comunidad, y qué le dejaré a quien viene detrás.
\end{softbox}

\textbf{Mi reflexión liberadora en cinco dimensiones.} Completa cada frase iniciada:

\small
\begin{tabularx}{\textwidth}{>{\bfseries\RaggedRight\arraybackslash}p{3.4cm}Y>{\RaggedRight\arraybackslash}p{4.6cm}}
\rowcolor{milcVino}\color{white}Dimensión & \color{white}\bfseries Mi respuesta (frame starter) & \color{white}\bfseries Algo que descubrí\\
1. Personal & Lo que más cambió en mí fue \linead{6.5cm} & \linead{4.2cm}\\
2. Emocional & La emoción más fuerte fue \linead{2.5cm} y la regulé \linead{3cm} & \linead{4.2cm}\\
3. Ciudadana & En democracia esto importa porque \linead{6cm} & \linead{4.2cm}\\
4. Local (Cartago) & En mi barrio sería útil para \linead{6.5cm} & \linead{4.2cm}\\
5. Intergeneracional & A mi abuela/abuelo le diría \linead{6.5cm} & \linead{4.2cm}\\
\end{tabularx}
\normalsize

\begin{infoband}{milcVino}
\textbf{Para la próxima sustentación.} Vuelve a esta tabla en seis meses. Verás que las respuestas cambian — no porque lo aprendido cambie, sino porque tú cambias. La evaluación liberadora es una conversación contigo a través del tiempo.
\end{infoband}


\puente{Y para terminar, tres voces sobre la hoja como base. Dussel sobre las categorías que invisibilizan, Marco Aurelio sobre la disciplina de la columna correcta, Floridi sobre la hoja como infraestructura mínima de cualquier organización digital.}

\begin{titlebox}{milcGradeDark}
Triángulo de pensamiento · cierre filosófico
\end{titlebox}

\begin{softbox}{milcVino}{milcGris}{Dussel · lente del nosotros}
\textit{``Una columna decide qué se cuenta como dato y qué queda fuera del modelo.''}

Cuando agregas la columna \emph{``ingreso''} pero no la columna \emph{``trabajo no remunerado''}, decides que el cuidado, la limpieza y la cocina familiar no son datos. Cuando tu base de datos de \emph{``clientes''} no tiene columna para \emph{``personas que no pueden pagar pero usan el servicio''}, esas personas dejan de existir para la administración. La phronesis del diseño de bases exige preguntar siempre: \emph{¿quién no tiene columna en mi hoja y por qué?}.

\textbf{En mi base de datos, ¿qué personas o situaciones quedaron sin columna? ¿Qué cambia en la decisión cuando agrego una columna para nombrarlos?}
\end{softbox}
\linead{\linewidth}\\[1mm]
\linead{\linewidth}

\begin{softbox}{milcMostaza}{milcGris}{Estoicismo · Marco Aurelio · cuidado interior}
\textit{``La columna correcta en su lugar es la base de toda obra que dura.''}

Es tentador improvisar la estructura y \emph{``ya arreglamos después''}. La sabiduría estoica recomienda lo contrario: \emph{la disciplina de la columna correcta al inicio ahorra horas después}. Una hoja bien estructurada desde la primera fila aguanta miles de filas; una mal estructurada se rompe al multiplicarse. La paciencia inicial es virtud profesional.

\textbf{¿Resistí la tentación de improvisar la estructura? ¿Qué ahorro de trabajo futuro gano con la disciplina inicial?}
\end{softbox}
\linead{\linewidth}\\[1mm]
\linead{\linewidth}

\begin{softbox}{milcTurquesa}{milcGris}{Floridi · lente de la infoesfera}
\textit{``La hoja de cálculo bien estructurada es la infraestructura mínima de cualquier organización contemporánea.''}

Antes de los sistemas ERP, antes de los CRM, antes de cualquier software de gestión, hubo y sigue habiendo hojas de cálculo. El 80\% de las microempresas de América Latina funciona con hojas. Aprender a estructurarlas como bases de datos a los 16 años es alfabetización profesional anticipada: el cimiento de cualquier transformación digital posterior.

\textbf{¿Cuántas decisiones organizacionales que vi esta semana descansaban en una hoja de cálculo bien o mal estructurada? ¿Qué cambio si dominio esa estructura desde ya?}
\end{softbox}
\linead{\linewidth}\\[1mm]
\linead{\linewidth}

\begin{titlebox}{milcVino}
Compromiso final
\end{titlebox}

\small
\begin{tabularx}{\textwidth}{>{\RaggedRight\arraybackslash}p{3.0cm}YYY}
\rowcolor{milcVino}\color{white}\bfseries Criterio & \color{white}\bfseries Lo logré & \color{white}\bfseries En proceso & \color{white}\bfseries Necesito apoyo\\
Comprensión del tema & Explico ideas centrales con mis palabras. & Reconozco ideas pero falta precisión. & Necesito repasar conceptos base.\\
Pensamiento computacional & Apliqué los 4 pilares al diseñar mi producto. & Apliqué algunos pilares. & Necesito releer la fase 2.\\
Aplicación y producto & Mi producto cumple los criterios y se puede revisar. & Producto funcional con ajustes pendientes. & Aún en construcción.\\
Reflexión 5 dimensiones & Respondí cada dimensión con honestidad. & Respondí algunas con detalle. & Mis respuestas son muy generales.\\
Triángulo de pensamiento & Conecté las 3 voces con mi trabajo real. & Conecté al menos una. & No relaciono las voces con mi proceso.\\
\end{tabularx}
\normalsize

\textbf{Hoy entendí que:}\\[1mm]
\linead{\linewidth}\\[1mm]
\linead{\linewidth}

\textbf{Mi compromiso para la próxima sesión es:}\\[1mm]
\linead{\linewidth}\\[1mm]
\linead{\linewidth}

\begin{softbox}{milcMostaza}{milcGris}{Cita de despedida · Marco Aurelio}
\textit{``El obstáculo en el camino se vuelve el camino. Lo que estorba al camino se vuelve camino.''}\\[1mm]
Cada error de hoy, bien observado, se convierte en pieza del camino que pisarás mañana.
\end{softbox}

\small
\begin{tabularx}{\textwidth}{>{\bfseries}p{3.6cm}Y}
\rowcolor{milcGradeDark!12}Nombre completo: & \linead{10cm}\\
Fecha: & \linead{4cm} \quad Curso/grupo: \linead{4cm}\\
Mi firma: & \linead{9cm}\\
\end{tabularx}
\normalsize

\begin{infoband}{milcGradeDark}
\textbf{Para la próxima semana.} Toma una hoja real de tu vida (registro de gastos, lista de contactos del grupo, notas del curso) y rediseñala según las 3 reglas tidy. Comparte la versión limpia con quien comparte la hoja contigo. La phronesis de la estructura se entrena en lo cotidiano, no solo en clase.
\end{infoband}

\end{document}
